% ============================================================================
%  SpeedrunGettingStarted.tex
%  CIS 3300 running example for Weeks 12 to 14: The Speedrun Leaderboard.
%  Getting Started: the introduction, how to use the guides, the data, and the starter.
%  Built from class-notes-template-cis3300.tex. Requires classnotes.sty in the
%  same folder. Compile twice with pdflatex so "Page X of Y" resolves.
% ============================================================================
\documentclass[11pt]{article}
\usepackage[margin=1in]{geometry}

\usepackage[pagenumbers]{classnotes}
% Latin Modern gives scalable fonts with a Unicode map, so Acrobat and screen
% readers can extract every character.
\usepackage{lmodern}

% Sets the PDF title and language. Both are accessibility requirements.
\classnotesmeta{CIS 3300 Guide: The Speedrun Leaderboard, Getting Started}

% Java is the default listing style; terminal output uses [style=sh], and wide
% program output uses [style=wide] so its lines fit without wrapping.
\lstset{style=java}
\lstdefinestyle{wide}{style=sh, basicstyle=\ttfamily\footnotesize}

\begin{document}
\sloppy

\classnotestitle{CIS 3300 \quad$\cdot$\quad THE SPEEDRUN LEADERBOARD}{Getting Started}{A running example, built one checkpoint at a time}

\noindent
Speedrunners compete to finish a video game as fast as possible, and a leaderboard records each
run: the player, the game and category, the time, the date, and whether moderators verified it. In
these guides we build a program that reads submitted runs from a file, rejects the bad ones, and
reports world records, personal bests, and rankings.

We build it in class, one checkpoint at a time, from Week 12 to Week 14, and these guides contain
every step, including any we do not reach. Each checkpoint prepares you for part of Lab W14, which
repeats the same process with weather data but different details, so you need to understand each
step, not copy it.

To see a real leaderboard, visit \url{https://www.speedrun.com/sm64}, which includes the 16 Star
and 120 Star categories used in our data. Our game titles and categories are real, but the players
and times are invented.

\notebox{A note about this guide.}{I developed this guide with help from AI, which drafted much of
the code, data, and explanations, and I review and revise them. Every program was compiled and
run, and every output shown is real. If something looks wrong, does not match what you see in
Eclipse, or contradicts the book or the notes, tell me. This guide shows how to build one program.
It does not replace the notes, the assigned chapters, or the recorded lectures.}

% ============================================================================
\section{How to Use These Guides}

Each checkpoint is a complete Eclipse project that compiles and runs, and each checkpoint has a
guide of its own. Every checkpoint guide starts from one project and ends at the next, as
Table~\ref{tab:projects} shows. That means you can start any checkpoint, even if you missed the
class before it: import the project its guide starts from, and work through the steps.

\begin{classtable}
\captionof{table}{The Eclipse projects, and where each one is used.}\label{tab:projects}
\renewcommand{\arraystretch}{1.3}
\begin{tabular}{>{\ttfamily}L{2.2in} L{3.8in}}
\toprule
\rowcolor{navy!8}
\textrm{\textbf{Project}} & \textbf{Use it to} \\
\midrule
SpeedrunStarter & Start the Starter Program, below, and Checkpoint 1 \\
SpeedrunCheckpoint1 & Check your Checkpoint 1 work, or start Checkpoint 2 \\
SpeedrunCheckpoint2 & Check your Checkpoint 2 work, or start Checkpoint 3 \\
SpeedrunCheckpoint3 & Check your Checkpoint 3 work, or start Checkpoint 4 \\
SpeedrunCheckpoint4 & Check your Checkpoint 4 work, or start Checkpoint 5 \\
SpeedrunCheckpoint5 & Check your Checkpoint 5 work, or start Checkpoint 6 \\
SpeedrunCheckpoint6 & Check your Checkpoint 6 work: the finished program \\
\bottomrule
\end{tabular}
\end{classtable}

\subsection{Importing a project into Eclipse}

\begin{enumerate}[leftmargin=1.6em]
  \item Download the project's zip file. You do not need to unzip it.
  \item In Eclipse, choose \textbf{File}, then \textbf{Import}. Open \textbf{General}, choose
        \textbf{Existing Projects into Workspace}, and click \textbf{Next}.
  \item Choose \textbf{Select archive file}, browse to the zip file, make sure the project is
        checked, and click \textbf{Finish}.
\end{enumerate}

Every project has the same layout: a \cmd{speedrun} package for the classes that represent the
data, an \cmd{app} package for the program with \cmd{main}, and a \cmd{files} folder for the
data files.

\subsection{How each checkpoint guide works}

Each checkpoint guide lists the project it starts from, the project it ends at, and the Lab W14 tasks it
prepares you for. Then it gives numbered steps. Each step shows the code to type and the output
you should see when you run it. If your output does not match, stop and compare your code with the
listing before going on, because each step builds on the one before it. Each guide ends with
\emph{Check yourself} questions and their answers.

% ============================================================================
\section{The Data}

The file \cmd{files/runs.csv} holds one submitted run per line, after a header line. The runs
are in the order they were submitted, not in any useful order:

\begin{lstlisting}[style=wide]
player,game,category,time,date,verified
NovaByte,Super Mario 64,16 Star,16:51.5,2025-09-03,no
Kestrel,The Legend of Zelda: Ocarina of Time,100%,4:05:54.7,2025-11-11,no
\end{lstlisting}

Times are written the way speedrunners write them: minutes, seconds, and tenths, such as
\cmd{16:51.5}, or hours first when a run takes an hour or more, such as \cmd{4:05:54.7}.

Not every line is good. Table~\ref{tab:rules} lists what the leaderboard accepts. A line that
does not have exactly six fields, or has a field the table does not accept, is a \emph{bad line}.
The file \cmd{files/runs_clean.csv} has eight good runs and no bad lines, for testing.

\begin{classtable}
\captionof{table}{What the leaderboard accepts in each field.}\label{tab:rules}
\renewcommand{\arraystretch}{1.3}
\begin{tabular}{>{\ttfamily}L{1.3in} L{4.7in}}
\toprule
\rowcolor{navy!8}
\textrm{\textbf{Field}} & \textbf{Accepted values} \\
\midrule
player, game, category & Text that is not blank \\
time & A time written as \cmd{M:SS.s} or \cmd{H:MM:SS.s}, with minutes and seconds below 60, and greater than 0 \\
date & A real date written as year, month, and day, such as \cmd{2026-02-14}, and not in the future \\
verified & Exactly \cmd{yes} or \cmd{no} \\
\bottomrule
\end{tabular}
\end{classtable}

% ============================================================================
\section{The Starter Program}

\textbf{Starts from:} SpeedrunStarter. \quad \textbf{Changes nothing:} this section only runs it.

The starter program reads the clean file, splits each line on commas, turns the time and the date
into Java objects, and prints each run. Here is all of it:

\begin{lstlisting}
public class SpeedrunReport {

    public static void main(String[] args) throws IOException {
        List<String> lines = Files.readAllLines(Path.of("files/runs_clean.csv"));

        // Line 0 is the header, so the runs start at line 1.
        for (int i = 1; i < lines.size(); i++) {
            String[] fields = lines.get(i).split(",");
            String player = fields[0];
            String game = fields[1];
            String category = fields[2];
            Duration time = RunTime.parse(fields[3]);
            LocalDate date = LocalDate.parse(fields[4]);
            boolean verified = fields[5].equals("yes");

            System.out.printf("%-10s %-36s %-8s %10s  %s  %s%n",
                    player, game, category, RunTime.format(time), date,
                    verified ? "verified" : "unverified");
        }
    }
}
\end{lstlisting}

Three things in it may be new:

\begin{itemize}[leftmargin=1.4em]
  \item \cmd{Files.readAllLines} reads the whole file into a \cmd{List<String>}, one element per
        line.
  \item \cmd{Duration} is Java's class for an amount of time, in the \cmd{java.time} package,
        next to \cmd{LocalDate}.
  \item \cmd{RunTime} is a small class in the \cmd{speedrun} package that is provided for you.
        \cmd{RunTime.parse} turns text such as \cmd{16:51.5} into a \cmd{Duration}, and
        \cmd{RunTime.format} turns a \cmd{Duration} back into text. Use it as it is.
\end{itemize}

\subsection{Step 1: Run it on the clean file}

Import SpeedrunStarter, open \cmd{SpeedrunReport.java} in the \cmd{app} package, and run it.

\begin{lstlisting}[style=wide]
NovaByte   Super Mario 64                       16 Star     16:51.5  2025-09-03  unverified
Tangent    Super Mario 64                       16 Star     15:42.6  2025-09-24  unverified
Kestrel    The Legend of Zelda: Ocarina of Time 100%      4:05:54.7  2025-11-11  unverified
Lumen      Super Mario 64                       16 Star     17:03.4  2025-11-15  verified
Lumen      Super Mario 64                       16 Star     17:06.4  2025-12-13  verified
NovaByte   Super Mario 64                       120 Star  1:43:33.5  2025-12-19  verified
Moss       The Legend of Zelda: Ocarina of Time 100%      3:50:02.6  2025-12-25  unverified
Sprocket   Super Mario 64                       120 Star  1:38:42.1  2025-12-28  verified
\end{lstlisting}

\subsection{Step 2: Run it on the real file}

Change \cmd{runs_clean.csv} to \cmd{runs.csv} in the file name, and run it again.

\begin{lstlisting}[style=wide]
NovaByte   Super Mario 64                       16 Star     16:51.5  2025-09-03  unverified
Tangent    Super Mario 64                       16 Star     15:42.6  2025-09-24  unverified
Kestrel    The Legend of Zelda: Ocarina of Time 100%      4:05:54.7  2025-11-11  unverified
Lumen      Super Mario 64                       16 Star     17:03.4  2025-11-15  verified
Lumen      Super Mario 64                       16 Star     17:06.4  2025-12-13  verified
Exception in thread "main" java.lang.NumberFormatException: For input string: "4x.1"
	at java.base/jdk.internal.math.FloatingDecimal.readJavaFormatString(FloatingDecimal.java:2054)
	at java.base/jdk.internal.math.FloatingDecimal.parseDouble(FloatingDecimal.java:110)
	at java.base/java.lang.Double.parseDouble(Double.java:792)
	at speedrun.RunTime.parse(RunTime.java:38)
	at app.SpeedrunReport.main(SpeedrunReport.java:30)
\end{lstlisting}

Read a stack trace from the top. The first line says \emph{what} went wrong: a
\cmd{NumberFormatException}, because \cmd{"4x.1"} is not a number. The lines below it say
\emph{where}, starting with the method that threw it and ending with \cmd{main}. Skip the lines
inside \cmd{java.base}, which are Java's own code, and look for the first line in code from this
project. Here, \cmd{RunTime.parse} threw the exception while \cmd{main} was calling it on line 30.

The bad line is line 7 of the file:

\begin{lstlisting}[style=wide]
Kestrel,Celeste,Any%,28:4x.1,2025-10-02,yes
\end{lstlisting}

Someone typed an \cmd{x} instead of a digit. The real problem is what happened next: nothing.
The program stopped, and the 36 good runs after that line were never read. One typo threw away
the whole leaderboard.

\subsection{Step 3: Run it with a missing file}

Change the file name to \cmd{run_clean.csv}, without the \cmd{s}, and run it again.

\begin{lstlisting}[style=wide]
Exception in thread "main" java.nio.file.NoSuchFileException: files/run_clean.csv
	at java.base/sun.nio.fs.UnixException.translateToIOException(UnixException.java:92)
	... (ten more lines inside java.base)
	at app.SpeedrunReport.main(SpeedrunReport.java:22)
\end{lstlisting}

The real trace has ten more lines of Java's own code where the \cmd{...} is; your line numbers
inside \cmd{java.base} may also differ, depending on your version of Java. The message itself is
useful, because it names the file, but a person running the program should never see a stack
trace. Change the file name back to \cmd{runs_clean.csv}.

\medskip\noindent\textbf{Next:} the Checkpoint 1 guide, which starts from SpeedrunStarter.

\end{document}
